%! Polynomial Functions and Rates of Change — Unit 2, Lesson 2.2 — Slides (Beamer)
\documentclass[aspectratio=169, 11pt]{beamer}
\usepackage{precalculus-beamer}   % provides \plumheader{} and \sectionlabel[color]{}

\begin{document}

% TITLE SLIDE
{
\setbeamercolor{background canvas}{bg=plum}
\begin{frame}[plain]
  \vspace{2.8cm}\hspace{0.55cm}%
  \begin{minipage}{0.9\textwidth}
    {\color{goldacc}\bfseries\small LESSON 2.2}\\[0.5cm]
    {\color{white}\bfseries\LARGE Polynomial Functions and Rates of Change}\\[1.2cm]
    {\color{lilacmid}\small Precalculus~$\cdot$~Unit 2: Polynomial Functions}
  \end{minipage}
\end{frame}
}

% HOOK SLIDE
\begin{frame}[t, plain]
  \plumheader{Hook: Two Summits, One Highest Point}
  \sectionlabel[goldacc]{Think About It}

  A trail guide lists two summits on the Ridge Trail:

  \bigskip
  \begin{center}
    {\LARGE \textbf{Bell Knob, 1{,}200 ft \qquad Signal Peak, 1{,}850 ft}}
  \end{center}

  \bigskip
  \begin{itemize}\setlength{\itemsep}{8pt}
    \item Are both of these really summits?
    \item Which one goes on the sign at the trailhead?
  \end{itemize}

  \bigskip
  \textit{Each is highest \textbf{nearby} --- a \textbf{local} maximum. Only
  Signal Peak is highest \textbf{anywhere} --- the \textbf{global} maximum.}
\end{frame}

% SECTION 1 — THE DEFINITION
\begin{frame}[t, plain]
  \plumheader{1. Polynomial Functions --- The Definition, Precisely}

  \begin{block}{Standard form}
    \[
      p(x) = a_n x^n + a_{n-1}x^{n-1} + \cdots + a_1 x + a_0
    \]
    where \(n\) is a positive integer, every \(a_i\) is real, and
    \(a_n \neq 0\).
  \end{block}

  \medskip
  \begin{itemize}\setlength{\itemsep}{8pt}
    \item \textbf{Degree:} \(n\) --- the greatest exponent on \(x\)
    \item \textbf{Leading term:} \(a_n x^n\) \quad
          \textbf{Leading coefficient:} \(a_n\)
    \item A constant function such as \(p(x) = 7\) has degree \textbf{0}
  \end{itemize}

  \bigskip
  \textit{Quick check:} \(f(x) = 4x^3 - x^2 + 9\) --- degree 3, leading term
  \(4x^3\), leading coefficient 4.
\end{frame}

% SECTION 2 — INCREASING, DECREASING, LOCAL EXTREMA
\begin{frame}[t, plain]
  \plumheader{2. Increasing, Decreasing, and Local Extrema}

  \begin{columns}[T]
    \column{0.48\textwidth}
      \begin{center}
      \begin{tikzpicture}[x=0.72cm, y=0.44cm]
        \draw[linegray, very thin, step=1] (0,0) grid (6,7);
        \draw[->, thick] (0,0) -- (6.6,0) node[below right] {\small \(x\)};
        \draw[->, thick] (0,0) -- (0,7.6) node[above] {\small \(f(x)\)};
        \foreach \x in {1,2,3,4,5,6} \node[below] at (\x,0) {\small \x};
        \foreach \y in {2,4,6} \node[left] at (0,\y) {\small \y};
        \draw[very thick, plum] plot [smooth] coordinates
          {(0,1) (0.8,3.3) (1.5,5.2) (2,6) (2.5,5.2) (3,4.0)
           (3.5,2.6) (4,2) (4.5,2.6) (5,3.6) (5.5,5.1) (6,7)};
        \fill[goldacc] (0,1) circle (2.5pt); \node[below right] at (0.05,1) {\small \(A\)};
        \fill[goldacc] (2,6) circle (2.5pt); \node[above] at (2,6.15) {\small \(B\)};
        \fill[goldacc] (4,2) circle (2.5pt); \node[below] at (4,1.9) {\small \(C\)};
        \fill[goldacc] (6,7) circle (2.5pt); \node[left] at (5.85,6.9) {\small \(D\)};
      \end{tikzpicture}
      \end{center}

    \column{0.48\textwidth}
      \begin{itemize}\setlength{\itemsep}{7pt}
        \item Increasing on \((0,2)\) and \((4,6)\); decreasing on \((2,4)\)
        \item \(B\): increasing \(\to\) decreasing, a \textbf{local maximum}
              --- value \textbf{6} at \(x=2\)
        \item \(C\): decreasing \(\to\) increasing, a \textbf{local minimum}
              --- value \textbf{2} at \(x=4\)
      \end{itemize}

      \medskip
      \begin{block}{Included endpoints count}
        The domain is \([0,6]\): \(A\) is a local \textbf{minimum} and \(D\) a
        local \textbf{maximum}, just for being included endpoints.
      \end{block}
  \end{columns}
\end{frame}

% SECTION 3 — LOCAL VS GLOBAL
\begin{frame}[t, plain]
  \plumheader{3. Local vs.\ Global Extrema}

  \begin{columns}[T]
    \column{0.48\textwidth}
      \begin{block}{Local maxima}
        \(7\) at \(D\), \quad \(6\) at \(B\)\\[4pt]
        Greatest of them \(\Rightarrow\)
        \textbf{global maximum \(= 7\)} at \(x = 6\)
      \end{block}

    \column{0.48\textwidth}
      \begin{block}{Local minima}
        \(1\) at \(A\), \quad \(2\) at \(C\)\\[4pt]
        Least of them \(\Rightarrow\)
        \textbf{global minimum \(= 1\)} at \(x = 0\)
      \end{block}
  \end{columns}

  \bigskip
  \begin{block}{The trap}
    \(B\) is a local maximum but \textbf{not} the global one --- the graph
    climbs higher later. ``Local'' means highest \textit{nearby}; ``global''
    means highest \textit{anywhere on the domain}.
  \end{block}

  \bigskip
  \textit{Say it carefully: a maximum is an \textbf{output value}; it
  \textbf{occurs at} an input value.}
\end{frame}

% SECTION 4 — WHAT ZEROS AND DEGREE GUARANTEE
\begin{frame}[t, plain]
  \plumheader{4. What the Zeros and the Degree Guarantee}

  \begin{columns}[T]
    \column{0.44\textwidth}
      \begin{center}
      \begin{tikzpicture}[x=0.68cm, y=0.36cm]
        \draw[linegray, very thin, step=1] (0,-3) grid (6,3);
        \draw[->, thick] (0,0) -- (6.6,0) node[below right] {\small \(x\)};
        \draw[->, thick] (0,-3.4) -- (0,3.6) node[above] {\small \(p(x)\)};
        \foreach \x in {1,2,3,4,5,6} \node[below left] at (\x,0) {\small \x};
        \draw[very thick, plum] plot [smooth] coordinates
          {(0.35,-2.7) (1,0) (1.5,1.5) (2,2) (2.5,1.5) (3,0)
           (3.5,-1.5) (4,-2) (4.5,-1.5) (5,0) (5.65,2.7)};
        \fill[plum] (1,0) circle (2.5pt);
        \fill[plum] (3,0) circle (2.5pt);
        \fill[plum] (5,0) circle (2.5pt);
      \end{tikzpicture}
      \end{center}

    \column{0.52\textwidth}
      \begin{block}{Between two zeros, a turn}
        Between every two distinct real zeros there is \textbf{at least one}
        local maximum or minimum --- the graph left the axis, so it has to
        turn to come back.
      \end{block}

      \medskip
      \(k\) distinct real zeros \(\Rightarrow\) at least \(k-1\) local extrema.
  \end{columns}

  \bigskip
  \begin{block}{Degree decides whether a global extremum exists}
    \textbf{Even degree:} has a global maximum \textit{or} a global minimum.
    \quad
    \textbf{Odd degree:} \textbf{neither} --- one end runs to \(+\infty\),
    the other to \(-\infty\).
  \end{block}
\end{frame}

% SECTION 5 — POINTS OF INFLECTION
\begin{frame}[t, plain]
  \plumheader{5. Points of Inflection --- Where the \textit{Rate} Turns}

  From Lesson 2.1: AROCs increasing \(\Rightarrow\) concave up; AROCs
  decreasing \(\Rightarrow\) concave down.

  \bigskip
  \begin{tabular}{|c|c|c|c|c|c|c|}
  \hline
  \(x\) & 0 & 1 & 2 & 3 & 4 & 5 \\
  \hline
  \(f(x)\) & 1 & 3 & 8 & 17 & 25 & 29 \\
  \hline
  \end{tabular}

  \bigskip
  AROCs over \([0,1], [1,2], [2,3], [3,4], [4,5]\): \quad
  \(2,\ 5,\ 9,\ 8,\ 4\)

  \bigskip
  \begin{itemize}\setlength{\itemsep}{7pt}
    \item The AROCs rise to 9, then fall --- the \textbf{rate} turns around
          near \(x = 3\).
    \item Concave \textbf{up} before \(x=3\), concave \textbf{down} after:
          \(x = 3\) is a \textbf{point of inflection}.
  \end{itemize}
\end{frame}

% THE CONTRAST
\begin{frame}[t, plain]
  \plumheader{Extremum vs.\ Inflection --- Keep Them Apart}

  \begin{columns}[T]
    \column{0.48\textwidth}
      \begin{block}{Local maximum or minimum}
        The \textbf{function} turns around.\\[6pt]
        The AROCs change \textbf{sign}
        (\(+\) to \(-\), or \(-\) to \(+\)).
      \end{block}

    \column{0.48\textwidth}
      \begin{block}{Point of inflection}
        The \textbf{rate of change} turns around.\\[6pt]
        The AROCs change \textbf{direction} --- they stop growing and start
        shrinking, without changing sign.
      \end{block}
  \end{columns}

  \bigskip
  \begin{block}{The example that settles it}
    In the table \(2,\ 5,\ 9,\ 8,\ 4\), \textbf{every AROC is positive}. The
    function is increasing the whole way, so it has \textbf{no local extrema}
    at all --- and it still has a point of inflection at \(x=3\).
  \end{block}
\end{frame}

% CLASS PRACTICE
\begin{frame}[t, plain]
  \plumheader{Class Practice}
  \sectionlabel[redacc]{Try It}

  \begin{enumerate}\setlength{\itemsep}{10pt}
    \item A polynomial \(g\) has exactly four distinct real zeros. At least
          how many local extrema must \(g\) have, and why?
    \item A polynomial \(h\) has AROCs \(4,\ 9,\ 15,\ 11,\ 6\) over five
          consecutive equal-length intervals covering \([0,5]\).
          \begin{enumerate}\setlength{\itemsep}{5pt}
            \item Does \(h\) have a local maximum or minimum on \([0,5]\)?
            \item Near which input is there a point of inflection?
          \end{enumerate}
    \item True or false: a polynomial of degree 7 must have a global
          minimum. Justify in one sentence.
  \end{enumerate}
\end{frame}

\end{document}
